% SY52210A
% Sauerstoffberieselung
% 14.0
% 2013-11-30
\label{m:sauerstoffberieselung}
Bei jedem Patienten, bei dem eine
lebensbedrohliche Störung einer vitalen Funktion eingetreten
ist oder einzutreten droht (\enquote{Notfallpatient}), soll,
sofern keine Kontraindiktaionen vorliegen, soviel Sauerstoff
verabreicht werden, sodass die Sauerstoffsättigung
(Sp\ce{O2}) im Bereich von \EE{\VonBis{94}{98}\PC*} erreicht wird.

\begin{enumerate}
  \item Kontraindikationen und Gegenanzeigen prüfen:
  \begin{itemize}
    \item COPD (\FullPrettyRef{AASS-topic:copd})
    \item Hyperventilationssyndrom, Hyperventilationstetanie
    (\FullPrettyRef{AASS-topic:hyperventailationssyndrom})
  \end{itemize}
  \item Situationsgerechte Dosierung je nach
  zugrundeliegender Erkrankung. Grundsätzlich soll ein
  Sp\ce{O2} von
\EE{\VonBis{94}{98}\PC*} erreicht werden. Steht keine Pulsoxymetrie zur Verfügung, ist als
Richtwert von einer Dosis von \EE{8\LPerMin*} auszugehen,
welche dem klinischen Zustand des Patienten angepasst werden
muss.
  \item Auswahl des Hilfsmittels:
  \begin{itemize}
    \item \E{Sauerstoff\-brille}
    \ce{O2}-Flow bis~5\LPerMin*,
    \item \E{Sauerstoff\-maske}
    \ce{O2}-Flow ab 5\LPerMin*,
    \item \E{Sauerstoff\-maske mit
    Reservoir} \ce{O2}-Flow ab 5\LPerMin*,
  \end{itemize}
  \item Patientenaufklärung
  \item Sauerstoffgerät einschalten und Durchflußrate
  einstellen
  \item Bei Verwendung einer Sauerstoff\-maske mit
  Reservoir: Reservoir füllen
  \item Sauerstoffbrille/-maske positionieren 
\end{enumerate} 

\cite{AsboeLehrmeinungRd2011-09}